% =====================================================================
% V-D. Generalización geográfica: Bolivia y Lima  (E4)
% =====================================================================
\subsection{Generalización geográfica: Bolivia y Lima}
\label{sec:generalizacion}

Evaluamos el mismo modelo, sin ningún ajuste, en dos eventos que no vio
(Tabla~\ref{tab:regiones}, Fig.~\ref{fig:brecha}). Definimos la brecha como la caída relativa
del IoU respecto a test, $(\mathrm{IoU}_{\text{test}}-\mathrm{IoU}_{\text{región}})/
\mathrm{IoU}_{\text{test}}$.

\begin{table}[t]
\centering
\caption{Generalización a eventos no vistos (U-Net: 3 semillas).}
\label{tab:regiones}
\begin{tabular}{lccc}
\toprule
Región & U-Net IoU & Otsu IoU & Brecha vs.\ test \\
\midrule
Test (referencia) & $0.669\pm0.003$ & 0.554 & --- \\
Bolivia           & $0.663\pm0.026$ & 0.718 & 0.8\,\% \\
Lima              & $0.003$         & 0.011 & 99.5\,\% \\
\bottomrule
\end{tabular}
\end{table}

\begin{figure}[t]
\centering
\includegraphics[width=\columnwidth]{f8_brecha_regiones.png}
\caption{IoU por región: test, Bolivia y Lima, frente a Otsu.}
\label{fig:brecha}
\end{figure}

\textbf{Bolivia: generaliza.} La brecha es menor que la variación entre semillas. La red es aquí
más conservadora que Otsu (precisión 0.883, omisión 0.272), y Otsu la supera porque Bolivia es
un caso favorable para un umbral: grandes extensiones de agua calma y oscura en una llanura de
inundación.

\textbf{Lima: colapsa.} El IoU cae prácticamente a cero, y también el de Otsu: el problema no es
la red, sino la señal. La Tabla~\ref{tab:medianas} lo muestra: en Lima, los píxeles que la
etiqueta llama inundados tienen la \emph{misma} retrodispersión que los no inundados (diferencia
$<1$\,dB), mientras que en train el agua es 5--7\,dB más oscura que el fondo. Ningún modelo que
mire una sola imagen puede separar dos clases con la misma distribución.

\begin{table}[t]
\centering
\caption{Medianas de retrodispersión (dB) por clase.}
\label{tab:medianas}
\begin{tabular}{llcc}
\toprule
Conjunto & Clase & VV & VH \\
\midrule
\multirow{2}{*}{Train} & agua     & $-14.92$ & $-23.00$ \\
                       & no agua  & $-9.58$  & $-16.03$ \\
\midrule
\multirow{2}{*}{Lima (17/03)} & inundado    & $-10.79$ & $-18.64$ \\
                              & no inundado & $-10.14$ & $-17.77$ \\
\bottomrule
\end{tabular}
\end{table}

Para entender por qué, descargamos dos escenas más de la misma órbita: antes (05/03) y después
(29/03) del evento. En la zona inundada, la mediana de VV pasa de $-9.63$ a $-8.75$ y vuelve a
$-9.63$\,dB, y la de VH de $-17.12$ a $-16.75$ y luego a $-17.37$\,dB, mientras una zona de
control permanece estable. Es decir, el día del evento la zona inundada se \emph{aclaró} en lugar
de oscurecerse. Con esto planteamos las siguientes hipótesis:
\begin{itemize}
  \item \textbf{H1 -- La etiqueta no describe lo que vio el radar.} La capa EMSN038 es la máxima
        extensión potencial, delineada con imágenes ópticas de días posteriores. Parte del área
        etiquetada pudo no tener agua visible a las 10:38 UTC del 17/03, o tener solo lodo ya
        depositado.
  \item \textbf{H2 -- Doble rebote y lodo saturado en zona urbana.} Entre edificaciones, el agua
        forma reflectores de esquina que \emph{aumentan} la señal~\cite{zhao2022urbanaware}, y
        el barro húmedo y los escombros de un huaico son rugosos, no especulares. Es consistente
        con el aclaramiento medido.
  \item \textbf{H3 -- Agua turbulenta en cauces angostos.} El flujo del huaico es rugoso, y los
        brazos inundados miden pocas decenas de metros: apenas unos píxeles de 10\,m, dominados
        por el ruido \emph{speckle}.
\end{itemize}
H2 se verifica de forma directa en la Sección~\ref{sec:fallos}. Una cuarta hipótesis, que la
etiqueta tratara el agua permanente de forma distinta a Sen1Floods11, se probó con JRC y se
descartó (Sección~\ref{sec:fallos}). Por todo esto reportamos Bolivia como el resultado
cuantitativo de generalización y Lima como un caso negativo diagnosticado. Lima muestra un
límite físico de la detección con intensidad de una sola fecha, más que un defecto del
entrenamiento.
